% ============================================================
% M2 换装轮副本树试迁 · 练习件 第1课时（§1.1.1 课1 空间向量的概念及线性运算）——qp-m3 换装版
% 源件：成卷/练习件/课时01/main.tex（只读对照基线，原树未动）；本件试迁作业仅在本副本树。
% 迁移（锚点迁移法，方案草案 §1）：答案册 16 键 练-课时01-1~16 逐题回嵌题后 ansblock；
%   源层锚＝体内行首 % ans:<键>；编译层锚＝\begin{ansblock}[键]（log M3-ANSKEY 两档恒发射）；
%   答案册两参 \ansline{解析}{…} 按 qp-m3 收编改名承入 \ansnote（qp-m3.sty:500 注）；
%   \anshang 双位档（题10起 7.6mm）照答案册 body.tex:678 同值迁入块内（组内局部）。
% 换装：六宏族 \input → \usepackage{qp-m3}；件型层收编宏（\huadianOverlapLx/\dangbadge/
%   \lxhang/\tihao/\lxopt/\xiexwei/\zu）删件内定义（qp-m3.sty:472-486 逐字同源收编）；
%   练习册档页脚三件（\qpxsyema/\lxshuziL/\lxshuziR）包未收编，件内保留。
% 回流：本件 \vbox＝0 处（grep 实证＝义务总表 §一.2③），课堂评价回流不适用；
%   尾页填充块 \tailfill 1 处（\end{multicols} 前，栏流内贴栏尾）。
% 模式：本文件＝true 档（含详解印本）；纯题档 main-false.tex 仅差 \usepackage[pure]{qp-m3}。
% 编译：xelatex main-true.tex（两遍）／xelatex main-false.tex（两遍）
% ============================================================
\documentclass[fontset=none]{ctexart}
\usepackage{qp-m3}

% ============================================================
% 练习册档页脚（件内保留件；\qpjianming/\qpzhangming 包内已有定义，renew）
% ============================================================
\renewcommand{\qpjianming}{练习件}
\renewcommand{\qpzhangming}{第一章\quad 空间向量与立体几何}
\renewcommand{\qpceming}{高中数学\quad 选择性必修第一册(人教B版)}% 包默认＝第二册(M3口径)，照原件 qp-headfoot.tex:24 拨回第一册——对勘[B3]实证必改项

% ---------- YJ-01 灰块页脚·练习册档（块 31.0×7.8mm，照练习线样张 31.0mm 修正版） ----------
\newcommand{\qpxsyema}{\ifnum\value{page}<10 00\fi\ifnum\value{page}>9 \ifnum\value{page}<100 0\fi\fi\arabic{page}}
\newcommand{\lxshuziL}{\hbox to 13.8mm{\setlength{\fboxsep}{0pt}\colorbox{pnumbg}{%
  \makebox[31.0mm][l]{\rule[-3.145mm]{0pt}{7.8mm}\hspace{2.8mm}\fontsize{10.7pt}{13pt}\selectfont\color{black}\numpnum\qpxsyema}}\hss}}
\newcommand{\lxshuziR}{\hbox to 13.8mm{\kern-17.2mm\setlength{\fboxsep}{0pt}\colorbox{pnumbg}{%
  \makebox[31.0mm][r]{\rule[-3.145mm]{0pt}{7.8mm}\fontsize{10.7pt}{13pt}\selectfont\color{black}\numpnum\qpxsyema\hspace{2.75mm}}}}}
\fancyfoot[RO]{\raisebox{-1.24mm}[0pt][0pt]{\qphuitext\qpzhangming\quad{\heiti\qpjianming}\hspace{3mm}\lxshuziL}}
\fancyfoot[LE]{\raisebox{-1.24mm}[0pt][0pt]{\lxshuziR\hspace{3mm}\qphuitext\qpceming}}

\begin{document}

% ============================================================
% 课时训练（三档 10＋4＋2＝16 题；答案 16 键逐题回嵌——\xiexwei 后插 ansblock，
% 书写位不废；短答灰底模（判据 §1.2.5：无 display 数学/表格/图、估高 <8 行））
% ============================================================
\keshi{第1课时\quad 空间向量的概念及线性运算}
\begin{multicols}{2}
\emergencystretch=1em

\dangbadge{基}{础}{巩}{固}

\tihao{1}\zu{概念·条件判定×选择} \tieside{简单(知识点三)}设\(\overrightarrow{a}\)，\(\overrightarrow{b}\)均为非零空间向量，则“\(\overrightarrow{a}\cdot\overrightarrow{b}=0\)”是“\(\overrightarrow{a}\perp\overrightarrow{b}\)”的\nobreak(\kongwei)
\lxopt{A．充分不必要条件；}
\lxopt{B．必要不充分条件；}
\lxopt{C．充要条件；}
\lxopt{D．既不充分也不必要条件}
\begin{ansblock}[练-课时01-1]
% ans:练-课时01-1
\ansitem{1}{C}
\end{ansblock}

\tihao{2}\zu{基底表示×填空} \tieside{简单(知识点二)}在平行六面体\(ABCD\text{-}\penalty100 A_1B_1C_1D_1\)中，\(\overrightarrow{AB}=\overrightarrow{a}\)，\(\overrightarrow{AD}=\overrightarrow{b}\)，\(\overrightarrow{AA_1}=\overrightarrow{c}\)，则\nobreak\(\overrightarrow{D_1B}=\)\kongbai{}．(用\(\overrightarrow{a}\)，\(\overrightarrow{b}\)，\(\overrightarrow{c}\)表示)
\begin{ansblock}[练-课时01-2]
% ans:练-课时01-2
\ansitem{2}{\(\overrightarrow{a}-\overrightarrow{b}-\overrightarrow{c}\)}
\end{ansblock}

\tihao{3}\zu{基底表示·定系数×填空} \tieside{简单(知识点二)}在平行六面体\(ABCD\text{-}\penalty100 A_1B_1C_1D_1\)中，点\(E\)在\(A_1C_1\)上，且\(\overrightarrow{A_1E}=\frac{1}{2}\overrightarrow{A_1C_1}\)，若\(\overrightarrow{AE}=x\overrightarrow{AA_1}+y(\overrightarrow{AB}+\overrightarrow{AD})\)，则\nobreak\(x=\)\kongbai{}，\(y=\)\kongbai{}．
\begin{ansblock}[练-课时01-3]
% ans:练-课时01-3
\ansitem{3}{\(x=1\)，\(y=\dfrac{1}{2}\)}
\end{ansblock}

\tihao{4}\zu{数量积运算×填空} \tieside{简单(知识点三)}已知\(|\overrightarrow{a}|=2\)，\(|\overrightarrow{b}|=3\)，\(\langle\overrightarrow{a},\overrightarrow{b}\rangle=60^\circ\)，则\nobreak\((\overrightarrow{a}+2\overrightarrow{b})\cdot(\overrightarrow{a}-\overrightarrow{b})=\)\kongbai{}．
\begin{ansblock}[练-课时01-4]
% ans:练-课时01-4
\ansitem{4}{\(-11\)}
\end{ansblock}

\tihao{5}\zu{数量积·夹角×填空} \tieside{中档(知识点三)}在长方体\(ABCD\text{-}\penalty100 A_1B_1C_1D_1\)中，\(AB=2\)，\(AD=2\)，\(AA_1=1\)，则向量\(\overrightarrow{AB_1}\)与\(\overrightarrow{AC}\)夹角的余弦值为\kongbai{}．
\begin{ansblock}[练-课时01-5]
% ans:练-课时01-5
\ansitem{5}{\(\dfrac{\sqrt{10}}{5}\)}
\ansnote{解析}{\(C\) 在底面！\(AB_{1}=(2,0,1)\)、\(AC=(2,2,0)\)，\(\cos=\dfrac{4}{\sqrt{5}\cdot 2\sqrt{2}}=\dfrac{\sqrt{10}}{5}\)。}
\end{ansblock}

\tihao{6}\zu{基底表示×填空} \tieside{简单(知识点二)}在直三棱柱\(ABC\text{-}\penalty100 A_{1}B_{1}C_{1}\)中，若\(\overrightarrow{CA} = \overrightarrow{a}\)，\(\overrightarrow{CB} = \overrightarrow{b}\)，\(\overrightarrow{CC_{1}} = \overrightarrow{c}\)，则\nobreak\(\overrightarrow{A_{1}B}\)=\kongbai{}.(用\(\overrightarrow{a},\overrightarrow{b},\overrightarrow{c}\)表示)
\begin{ansblock}[练-课时01-6]
% ans:练-课时01-6
\ansitem{6}{\(\overrightarrow{b}-\overrightarrow{a}-\overrightarrow{c}\)}
\end{ansblock}

\tihao{7}\zu{共面判定×解答} \tieside{简单(知识点三)}已知\(\overrightarrow{MP}=3\overrightarrow{MA}-2\overrightarrow{MB}\)，判断\(P\)，\(M\)，\(A\)，\(B\)四点是否共面，并说明理由．
\xiexwei{20mm}
\begin{ansblock}[练-课时01-7]
% ans:练-课时01-7
\ansitem{7}{共面}
\ansnote{解析}{\(\overrightarrow{MP}=3\overrightarrow{MA}-2\overrightarrow{MB}\) 系数和 \(3+(-2)=1\)。}
\end{ansblock}

\tihao{8}\zu{数量积·距离×填空} \tieside{简单(知识点三)}在大小为\(60^\circ\)的二面角\(A\text{-}\penalty100 EF\text{-}\penalty100 D\)中，四边形\(ABFE\)，\(CDEF\)都是边长为\(1\)的正方形，则\nobreak\(B\)，\(D\)两点间的距离为\kongbai{}．
\begin{ansblock}[练-课时01-8]
% ans:练-课时01-8
\ansitem{8}{\(\sqrt{2}\)}
\end{ansblock}

\tihao{9}\duoxuan\zu{数量积辨析多选×多选} \tieside{简单(知识点三)}设\(\overrightarrow{a}\)，\(\overrightarrow{b}\)为空间中的任意两个非零向量，下列各式中正确的有(\kongwei)
\lxopt{A．\({\overrightarrow{a}}^{2} = | \overrightarrow{a} |^{2}\)；}
\lxopt{B．\(\frac{\overrightarrow{a} \cdot \overrightarrow{b}}{\overrightarrow{a} \cdot \overrightarrow{a}} = \frac{\overrightarrow{b}}{\overrightarrow{a}}\)；}
\lxopt{C．\(( \overrightarrow{a} \cdot \overrightarrow{b} )^{2} = {\overrightarrow{a}}^{2} \cdot {\overrightarrow{b}}^{2}\)；}
\lxopt{D．\(( \overrightarrow{a} - \overrightarrow{b} )^{2} = {\overrightarrow{a}}^{2} - 2\overrightarrow{a} \cdot \overrightarrow{b} + {\overrightarrow{b}}^{2}\)}
\begin{ansblock}[练-课时01-9]
% ans:练-课时01-9
\ansitem{9}{AD}
\end{ansblock}

\tihao{10}\zu{翻折·距离×填空} \tieside{中档(知识点三)}已知矩形\(ABCD\)中，\(AB=1\)，\(BC=\sqrt{3}\)，将矩形\(ABCD\)沿对角线\(AC\)折起，使二面角\(B\text{-}\penalty100 AC\text{-}\penalty100 D\)的大小为\(60^\circ\)，则\nobreak\(|\overrightarrow{BD}|=\)\kongbai{}．
\begin{ansblock}[练-课时01-10]
% ans:练-课时01-10
\setlength{\anshang}{7.6mm}% 题10起双位档（照答案册 body.tex:678 同值，块内局部）
\ansitem{10}{\(\dfrac{\sqrt{7}}{2}\)}
\end{ansblock}

\dangbadge{综}{合}{提}{升}

\tihao{11}\zu{W7 空间四边形中点链化简×计算题组} \tieside{简单}在空间四边形\(ABCD\)中，连接\(AC\)，\(BD\)，设\(M\)，\(G\)分别是\(BC\)，\(CD\)的中点，化简下列各向量表达式：
\lxopt{(1)\(\overrightarrow{AB}+\overrightarrow{BC}+\overrightarrow{AD}\)；}
\lxopt{(2)\(\overrightarrow{AD}-\frac{1}{2}(\overrightarrow{AB}+\overrightarrow{AC})\)．}
\xiexwei{44mm}
\begin{ansblock}[练-课时01-11]
% ans:练-课时01-11
\ansitem{11}{(1)\(2\overrightarrow{AG}\)；(2)\(\overrightarrow{MD}\)}
\end{ansblock}

\tihao{12}\zu{W17 正方体夹角辨析选型×单选} \tieside{简单}在正方体\(ABCD\text{-}\penalty100 A_1B_1C_1D_1\)中，下列各对向量的夹角为\(45^\circ\)的是\nobreak(\kongwei)
\lxopt{\makebox[\dimexpr0.5\linewidth-0.5\lxhang-0.25em\relax][l]{A．\(\overrightarrow{AB}\)与\(\overrightarrow{A_1C_1}\)；}\makebox[\dimexpr0.5\linewidth-0.5\lxhang-0.25em\relax][l]{B．\(\overrightarrow{AB}\)与\(\overrightarrow{C_1A_1}\)；}}
\lxopt{\makebox[\dimexpr0.5\linewidth-0.5\lxhang-0.25em\relax][l]{C．\(\overrightarrow{AB}\)与\(\overrightarrow{A_1D_1}\)；}\makebox[\dimexpr0.5\linewidth-0.5\lxhang-0.25em\relax][l]{D．\(\overrightarrow{AB}\)与\(\overrightarrow{B_1A_1}\)}}
\begin{ansblock}[练-课时01-12]
% ans:练-课时01-12
\ansitem{12}{A}
\ansnote{解析}{\(AB\) 与 \(A_{1}C_{1}\) 成 \(45^{\circ}\)；B \(135^{\circ}\)、C \(90^{\circ}\)、D \(180^{\circ}\) 逐项坐标验。}
\end{ansblock}

\tihao{13}\zu{W22 正方体中点向量表示×解答} \tieside{中档}在正方体\(ABCD\text{-}\penalty100 A_1B_1C_1D_1\)中，点\(M\)，\(N\)分别是面对角线\(A_1B\)与\(B_1D_1\)的中点，若\(\overrightarrow{DA}=\overrightarrow{a}\)，\(\overrightarrow{DC}=\overrightarrow{b}\)，\(\overrightarrow{DD_1}=\overrightarrow{c}\)，将\(\overrightarrow{MN}\)用向量\(\overrightarrow{a}\)，\(\overrightarrow{b}\)，\(\overrightarrow{c}\)表示出来．
\xiexwei{80mm}
\begin{ansblock}[练-课时01-13]
% ans:练-课时01-13
\ansitem{13}{\(\overrightarrow{MN}=-\dfrac{1}{2}\overrightarrow{a}+\dfrac{1}{2}\overrightarrow{c}\)}
\ansnote{解析}{\(=\dfrac{1}{2}(\overrightarrow{BB_{1}}+\overrightarrow{A_{1}D_{1}})\)。}
\end{ansblock}

\tihao{14}\zu{W22 中位线平行且等于\(\frac{1}{2}\)BD 证明} \tieside{中档}已知正方体\(ABCD\text{-}\penalty100 A_1B_1C_1D_1\)中，\(M\)与\(N\)分别是棱\(BB_1\)与对角线\(CA_1\)的中点．求证：\(MN\parallel BD\)且\(MN=\frac{1}{2}BD\)．
\xiexwei{80mm}
\begin{ansblock}[练-课时01-14]
% ans:练-课时01-14
\ansitem{14}{证明见解析}
\ansnote{解析}{\(MN\parallel BD\) 且 \(MN=\dfrac{1}{2}BD\)（中位线）。}
\end{ansblock}

\dangbadge{思}{维}{探}{索}

\tihao{15}\zu{W24 正四面体基底数量积×六小问题组} \tieside{难}已知四面体\(ABCD\)的每条棱长都等于\(a\)，点\(E\)，\(F\)，\(G\)分别是棱\(AB\)，\(AD\)，\(DC\)的中点，求下列向量的数量积：
\lxopt{\makebox[\dimexpr0.5\linewidth-0.5\lxhang-0.25em\relax][l]{(1)\(\overrightarrow{AB}\cdot\overrightarrow{AC}\)；}\makebox[\dimexpr0.5\linewidth-0.5\lxhang-0.25em\relax][l]{(2)\(\overrightarrow{AD}\cdot\overrightarrow{DB}\)；}}
\lxopt{\makebox[\dimexpr0.5\linewidth-0.5\lxhang-0.25em\relax][l]{(3)\(\overrightarrow{GF}\cdot\overrightarrow{AC}\)；}\makebox[\dimexpr0.5\linewidth-0.5\lxhang-0.25em\relax][l]{(4)\(\overrightarrow{EF}\cdot\overrightarrow{BC}\)；}}
\lxopt{\makebox[\dimexpr0.5\linewidth-0.5\lxhang-0.25em\relax][l]{(5)\(\overrightarrow{FG}\cdot\overrightarrow{BA}\)；}\makebox[\dimexpr0.5\linewidth-0.5\lxhang-0.25em\relax][l]{(6)\(\overrightarrow{GE}\cdot\overrightarrow{GF}\)．}}
\xiexwei{100mm}
\begin{ansblock}[练-课时01-15]
% ans:练-课时01-15
\ansitem{15}{(1)\(\dfrac{a^{2}}{2}\)；(2)\(-\dfrac{a^{2}}{2}\)；(3)\(-\dfrac{a^{2}}{2}\)；(4)\(\dfrac{a^{2}}{4}\)；(5)\(-\dfrac{a^{2}}{4}\)；(6)\(\dfrac{a^{2}}{4}\)}
\end{ansblock}

\tihao{16}\zu{W25 面内动点数量积最值 PA·(PB+PC)×解答} \tieside{难}已知\(\triangle ABC\)是边长为\(2\)的等边三角形，\(P\)为平面\(ABC\)内的一点，求\(\overrightarrow{PA}\cdot(\overrightarrow{PB}+\overrightarrow{PC})\)的最小值．
\xiexwei{100mm}
\begin{ansblock}[练-课时01-16]
% ans:练-课时01-16
\ansitem{16}{最小值 \(-\dfrac{3}{2}\)；\(P\) 为高线 \(AM\) 中点}
\ansnote{解析}{\(f=2\left[x^{2}+y^{2}-\sqrt{3}y\right]\) 配方。}
\end{ansblock}

\tailfill

\end{multicols}
\end{document}
